% ======================= 前 置 页 =======================
\clearpage
\makeatletter\@mainmatterfalse\makeatother
\pagenumbering{roman}
\pagestyle{frontmatterpage}

% ---------------- 简装扉页 ----------------
\thispagestyle{empty}
\begingroup
\centering
\vspace*{7mm}

{\sffamily\footnotesize\color{pencil}
《远东异端史料与文明起源重构丛书》\ 第\,VII\,卷\par}

\vspace{27mm}

{\fontsize{26pt}{34pt}\selectfont\bfseries 战争的雷霆\par}
\vspace{8mm}
{\large ——记一种可能的中国历史——\par}

\vfill

{\sffamily\footnotesize\color{pencil}
МОСКОВСКИЙ УНИВЕРСИТЕТ\ \ ·\ \ САНКТ-ПЕТЕРБУРГ\ \ ·\ \ ФОНД «ВФС»\par}
\vspace{4mm}
{\footnotesize\color{pencil} 2024\par}
\endgroup
\clearpage

\seriespage{vii}

% ---------------- 正式扉页 ----------------
\thispagestyle{empty}
\begingroup
\centering
\vspace*{3mm}

{\sffamily\footnotesize\color{pencil}
《远东异端史料与文明起源重构丛书》\ 第\,VII\,卷\par}

\vspace{17mm}
{\fontsize{31pt}{40pt}\selectfont\bfseries 战争的雷霆\par}
\vspace{6mm}
{\LARGE 记一种可能的中国历史\par}

\vspace{10mm}
{\color{rulegrey}\rule{46mm}{0.4pt}}
\vspace{10mm}

{\itshape\small\color{graphite}
THUNDER OF WAR\\[2pt]
A Chronological Reconstruction of a Possible Chinese History\par}

\vspace{18mm}

{\large ［俄］А.\,В.\,巴兰尼科夫\ \ 著\par}
\vspace{2mm}
{\small\color{graphite} 莫斯科国立大学历史学副博士\par}
\vspace{1mm}
{\small\color{graphite} 远东文明重构课题组组长\par}

\vspace{11mm}

{\large 李\quad 文\quad 亚\qquad 张\quad 建\quad 国\ \ 译\par}
\vspace{2mm}
{\small\color{graphite} 中国社会科学院历史研究所伪史学研究中心\par}

\vfill

{\small\sffamily 莫斯科大学出版社\par}
{\small\sffamily 圣彼得堡学术书局\par}
{\small\sffamily VFS 基金会\par}
\vspace{3mm}
{\footnotesize 2024\par}
\endgroup
\clearpage

% ---------------- 版权页 ----------------
\thispagestyle{empty}
\begingroup
\footnotesize
\linespread{1.22}\selectfont
\centering
\begin{minipage}[t][\dimexpr\textheight-24pt\relax][s]{\textwidth}
\centering

{\sffamily\scriptsize\color{pencil}
МОСКОВСКИЙ ГОСУДАРСТВЕННЫЙ УНИВЕРСИТЕТ им.\ М.\,В.\,ЛОМОНОСОВА\\
ИЗДАТЕЛЬСТВО\ \ ·\ \ САНКТ-ПЕТЕРБУРГСКОЕ АКАДЕМИЧЕСКОЕ ИЗДАТЕЛЬСТВО\\
ФОНД «ВФС»\ \ ·\ \ VFS FOUNDATION\par}

\vspace{3.2mm}
{\bfseries 图书在版编目（CIP）数据\par}
\vspace{2mm}

\begin{minipage}{0.94\textwidth}
\raggedright
战争的雷霆：记一种可能的中国历史／（俄罗斯）阿列克谢·瓦西里耶维奇·巴兰尼科夫著；
李文亚，张建国译．— 4版．据俄文第 3 版译出．— 莫斯科：莫斯科大学出版社；圣彼得堡：圣彼得堡学术书局；VFS 基金会，2024.3
（远东异端史料与文明起源重构丛书；第\,VII\,卷）
ISBN 978-5-317-07423-4

\vspace{1.4mm}
\hspace{2em}Ⅰ．①战…\quad Ⅱ．①巴…②李…③张…\quad Ⅲ．①中国历史－重构史学－研究\ \ Ⅳ．①K207.8

\vspace{1.4mm}
\hspace{2em}中国国家版本馆 CIP 数据核字（2024）第 0316 号
\end{minipage}

\vspace{3.2mm}
{\color{rulegrey}\rule{0.94\textwidth}{0.3pt}}
\vspace{2.4mm}

\begin{minipage}{0.94\textwidth}
\raggedright
\renewcommand{\arraystretch}{1.0}
\begin{tabular}{@{}p{15mm}p{0.76\textwidth}@{}}
书\quad 名 & 战争的雷霆——记一种可能的中国历史\\
著\quad 者 & ［俄］А.\,В.\,巴兰尼科夫\\
译\quad 者 & 李文亚\quad 张建国\\
责任编辑 & Н.\,П.\,谢列兹尼奥娃\quad （中文版）王守拙\\
特约编辑 & А.\,И.\,库兹涅佐夫\\
封面设计 & С.\,М.\,别列津\\
版式设计 & 圣彼得堡学术书局版式室\\
责任校对 & 赵无咎\\
出版发行 & 莫斯科大学出版社\quad 圣彼得堡学术书局\quad VFS 基金会\\
资料提供 & VFS 基金会异常文献部（控制·收容·保护）\\
社\quad 址 & 莫斯科市列宁山，出版大楼（119991）\\
网\quad 址 & www.msu-press.ru\\
电子信箱 & barannikov@fareast-recon.msu.ru\\
印\quad 刷 & 圣彼得堡第一学术印务厂\\
开\quad 本 & 185\,mm × 260\,mm\quad 1/16\\
印\quad 张 & 4.13\\
字\quad 数 & 73 千字\\
版\quad 次 & 2024 年 3 月第 4 版\\
印\quad 次 & 2024 年 3 月第 1 次印刷\\
印\quad 数 & 2501—4000 册\\
书\quad 号 & ISBN 978-5-317-07423-4\\
定\quad 价 & 128.00 元\\
\end{tabular}
\end{minipage}

\vspace{3mm}

\begin{minipage}{0.94\textwidth}
\raggedright
{\scriptsize\color{pencil}
版权所有\ \ 侵权必究。未经出版者书面许可，不得以任何方式复制或抄袭本书之部分或全部内容。\\
本书如有印装质量问题，请与本社读者服务部联系调换。\\
本书采用无酸纸印刷，符合 ISO 9706 文献耐久性标准。\par}
\end{minipage}

\vfill

% ---- 条形码（真实 EAN-13 图案）----
\begin{tikzpicture}[baseline=0pt]
  \draw[ink, line width=0.3pt] (0,0) rectangle (44mm,16mm);
  % 模块宽 0.366mm × 95 = 34.77mm；起点 4.6mm，左右留白对称
  \begin{scope}[shift={(4.6mm,5.4mm)}]
    \EANthirteen{0.366mm}{7.4mm}
  \end{scope}
  % 人读字符：首位居左静区，左六位与右六位分别居中于对应半幅
  \node[font=\sffamily\fontsize{4.4pt}{5pt}\selectfont, anchor=base]
    at (3.1mm,3.9mm) {9};
  \node[font=\sffamily\fontsize{4.4pt}{5pt}\selectfont, anchor=base]
    at (13.4mm,3.9mm) {785317};
  \node[font=\sffamily\fontsize{4.4pt}{5pt}\selectfont, anchor=base]
    at (30.6mm,3.9mm) {074234};
  \node[font=\sffamily\fontsize{3.8pt}{4.4pt}\selectfont, anchor=base]
    at (22mm,14.6mm) {ISBN 978-5-317-07423-4};
\end{tikzpicture}

\end{minipage}
\endgroup
\clearpage

% ---------------- 献词页 ----------------
\thispagestyle{empty}
\vspace*{58mm}
\begin{flushright}
{\itshape\large\color{graphite}
献给我的同行们——\\
那些尚未被“生态重启”抹去的人\par}
\vspace{7mm}
{\small\color{pencil}\itshape А.\,В.\,Б.\par}
\end{flushright}
\vfill
\clearpage

\pagestyle{fancy}

% ======================= 译 者 序 =======================
\clearpage
\chapter*{译\ 者\ 序}
\markboth{译者序}{译者序}
\addcontentsline{toc}{chapter}{译者序}

\noindent
二〇二一年十二月，莫斯科，列宁山。气温零下二十七度，莫斯科大学主楼十七层的供暖管道发出类似风洞实验的嗡鸣。巴兰尼科夫教授从保险柜里取出一只锈迹斑斑的铁盒，盒盖上的火漆早已开裂，编号仍清晰可辨：\textsc{ф.\,117, оп.\,4, д.\,2187}。他说：“这是您在中国国内永远看不到的东西。”\footnote{苏联科学院远东委员会（1949—1991）的档案于一九九二年部分解密。第 117 全宗第 4 目第 2187 号卷宗至今未见于任何公开目录，其存在仅由两份交叉引用的移交清单间接证明。参见俄罗斯科学院档案馆编《解密档案目录补编（1992—2010）》，莫斯科，2011 年，第 288 页。}

盒内是三十七份打字稿、四卷缩微胶片，以及一枚直径不足两厘米的金属圆片。圆片表面刻有极细的同心圆纹，肉眼几乎不可见。我最初的判断是“伪史料”——这类东西我们在潘家园见得太多。三个月后，莫斯科大学地质系对该圆片作了铼—锇同位素稀释质谱分析，其 \textsuperscript{187}\textsc{re}/\textsuperscript{187}\textsc{os} 比值为 3.147±0.006，与阿波罗十七号所采月海高地玄武岩的测定值在误差范围内完全一致。\footnote{分析报告编号 МГУ-ГЕО-2022-014。报告结论部分有一句被作者划掉的话：“建议不再重复采样。”}

我必须承认，那一刻我的学术信念出现了裂缝。而这本书，正是那条裂缝的完整剖面图。

\section*{一、\ 铁盒里的圆片}

巴兰尼科夫教授的治学路径，与我们所熟悉的任何一派中国史研究都不同。他不做考据，虽然他通读《旧唐书》《明实录》；他不做辨伪，虽然他一生都在辨伪。他所做的是“重构”——从一个文明留下的全部物质与文本痕迹中，反推它曾经拥有、后来又被迫遗忘的技术条件。

用他自己的话说：\kw{“历史学家不该问‘这不可能’，而该问‘要让它可能，需要多少能量’。”}\footnote{巴兰尼科夫《重构史学导论》，莫斯科大学出版社，2015 年，第 9 页。作者在此页页边手书：“此句被引用的次数，远超我被阅读的次数。”}

这一方法论被他概括为“四重证据法”：文献、器物、同位素、轨道。前三重，中国学者并不陌生，王国维的“二重证据法”已开其端；第四重则是他的独门。他坚持认为，任何大规模技术活动都会在近地轨道留下不可磨灭的残骸层，而这一层“轨道考古学”的档案，比任何正史都更难篡改。\footnote{据欧洲空间局 2019 年发布的轨道碎片统计，地球轨道上可追踪的十厘米以上人造物体约三万四千件，其中约百分之零点三“来源不明”。巴兰尼科夫在全书各处反复引用这百分之零点三，并称之为“我们唯一的诚实证人”。}

正是这一方法，使他在“莫斯科档案馆未公开的冷战时期考古卷宗”中，读出了中国学界从未设想过的东西：一个被反复重启的文明，一段被系统删除的技术史，以及一场在正史中仅占两行字的、星际规模的战争。

\section*{二、\ 四重证据法与它的裂缝}

坦率地说，本书的核心结论是惊世骇俗的，而我不打算在此为它辩护。

一个中国历史学者读到第一章时，大概会像我初见那枚圆片时一样，先笑，再怒，然后沉默。笑，是因为“月球广寒宫基地”这样的说法实在太像科幻小说；怒，是因为作者对《尚书》《史记》的处理方式近乎粗暴；沉默，则是因为——各位读者很快会发现——他引用的那些材料，页码是对的，字句是对的，版本是对的，甚至连避讳缺笔都一丝不差。

这是本书最令人不安之处：\kw{它几乎全部建立在真实的史料之上，只是把史料之间的空白，用一种我们无法接受、却也难以证伪的方式填满了。}\footnote{关于这种“以真实材料支撑不可证结论”的写法，可参见张建国《伪史学的技术伦理》，载《史学理论研究》2022 年第 3 期，第 55—71 页。该文认为此类写作“在方法上无可指责，在结论上不可接受”，并建议“将其归入文学”。}

中国社会科学院历史研究所伪史学研究中心之所以决定承担此书的翻译，理由正在于此：一种谬误，若其每一个步骤都合乎规范，那么它就不再是简单的谬误，而是一份关于“规范本身”的诊断书。我们译它，不是因为我们相信它，而是因为我们想知道，一种如此严整的荒谬，究竟是怎样被建造出来的。

\section*{三、\ 翻译体例说明}

翻译过程中的技术性说明，列于下：

其一，全书俄文原注保留原有格式，中文补注以“译者注”标明，不作合并。原注中大量使用的缩略语（如 \textsc{цга рф}、\textsc{ив ран}）按俄罗斯档案学界通例译出，首次出现时附原文。

其二，作者惯用的一组术语，中译取其字面义而非通行义，例如 \ru{перезапуск} 译作“重启”而不译作“复兴”，\ru{поле забвения} 译作“记忆消除场”而不译作“遗忘场”。这一处理曾与作者本人通信商定，他在回信中说：“随您怎么译。反正它们都会被删掉。”\footnote{作者 2022 年 3 月 14 日致译者电子邮件，现存译者处。此句之后尚有一行，因涉及未刊史料，遵作者嘱暂不发表。}

其三，本书第三版较前两版增补了第一章第三节的全部内容，并恢复了 1979 年苏联 censorship 删去的关于“大气适应性退化”的两页半文字。据作者称，这两页半是“全书的钥匙”。\footnote{删节本的影印件见巴兰尼科夫《被删去的两页半》，载《远东问题》1993 年第 2 期。该刊同期编者按称：“本刊发表此文，不代表本刊同意其结论。”}

其四，书中所有插图，均为作者据档案原件重绘。原件的照片因“技术原因”未能随书出版，这一点我们已多次要求说明，未获答复。

\section*{四、\ 一个译者的梦}

最后，请允许我以一个译者的身份，说一句不太像学者的话。

翻译这本书的十四个月里，我三次梦到长安。第一次，我站在含元殿的丹墀上，看见殿后有一座巨大的银色圆塔，塔身有水汽升腾；第二次，我站在萨尔浒的山谷里，天上有一个太阳，地上有十三个；第三次，我什么都没看见，只听见雷声。

醒来时我总在想：如果巴兰尼科夫是对的，那么我们的史书，究竟是一份记录，还是一份被精心保管的、给幸存者的说明书？

这个问题的答案，不在序言里，在正文中。

\vspace{12mm}
\begin{flushright}
{\small 李\ 文\ 亚}\\[2pt]
{\footnotesize\color{pencil} 二〇二三年九月，于北京建国门内大街五号}
\end{flushright}

% ======================= 中文版序言 =======================
\clearpage
\chapter*{中文版序言}
\markboth{中文版序言}{中文版序言}
\addcontentsline{toc}{chapter}{中文版序言}

\noindent
本书的中译者在序言中写道，我的结论“惊世骇俗”。这一说法客气得近乎虚伪。更准确的表述是：我的结论之所以令人不适，并非因为它们违背了史料，而是因为它们违背了诸位对史料的期待。

中国的史料是世界上最完备的史料，这一点我从不否认，并且由衷敬佩。但也正因如此，中国的史学传统形成了一种可以称为文本崇拜的习惯：凡不见于文字者，即视为不存在；凡见于文字者，即视为曾经发生。这是一种极为经济的世界观，它把整个宇宙压缩进一部可以随身携带的丛书之中。\footnote{作者此处原有一长注，列举中国类书对异域记载的系统性简化，后整段删除。据其助手称，删去的理由是“太刻薄了，而且他们译不过来”。}

各位即将读到的四章，分别处理四个时间断面。它们并非四个孤立的奇谈，而是同一条曲线的四个采样点，即一个文明如何获得技术，如何滥用技术，如何被封存，以及如何被自己的封存所出卖。

在此预先声明三点。

第一，本书不提供证据链，只提供证据点。点与点之间的连线是我的判断，凡判断皆可能有误。我欢迎反驳，但请反驳连线，不必反驳那些点；那些点是我用三十一年时间从十一个国家的档案中逐一辑出的。\footnote{作者此处用了一个俄语俗语 \ru{до мозолей}（直到磨出茧子），译者按字面译出。}

第二，本书不承认民族感情可以作为学术标准。若某一结论令某一民族感到不快，那只能说明历史本身曾经令人不快，而不能说明结论有误。历史学不是安慰剂。

第三，也是最重要的一点：本书第三版恢复了 1979 年被删去的两页半。那两页半讨论的是一个技术细节，即大气适应性退化对减速装置的影响系数。这一系数，以最朴素的方式证明了：我们的祖先不是从树上下来的，是从天上掉下来的。

这句话，我在莫斯科讲了四十三年。今天，它第一次被译成中文。

愿它被读完，而不只是被引用。

\vspace{10mm}
\begin{flushright}
{\small А.\,В.\,巴兰尼科夫}\\[2pt]
{\footnotesize\color{pencil} 二〇二三年十一月，于莫斯科列宁山}
\end{flushright}


% ---------------- 第四版说明 ----------------
\clearpage
\chapter*{第四版说明}
\markboth{第四版说明}{第四版说明}
\addcontentsline{toc}{chapter}{第四版说明}
\vspace*{-7mm}

这一版不在计划之内。

去年春天校读清样，我在第二章读到一段话，写得比我记得的好。我不记得写过它。

起初以为是排版时补的。查了原稿，是我的字。

我这一行有个笨规矩：凡自己写过却认不出的东西，不要删，要留着。留着，下次重读还会认不出。如果两次认不出的地方不在同一处，那就说明不是记错了\nobreak。

于是我照着做。此后几个月，把此前的几个本子逐页重读了一遍，把每一处对不上的地方标下来。

标记里属于原文的那些，改在了俄文第三版；二〇二三年九月印行。

俄文第三版改了四处。

第一处是第一章第四节。旧稿把仰韶、龙山、二里头的聚落说成降落之后的居民所造，与遗址年代不合，最晚的二里头也比撤离早一千三百年。这一版改由半坡、二里头的发掘材料入手，并补了一段限定：本章开篇所说的“早期形态”指技术与制度，不指人口。

第二处是第二章的四神一节。那一节在俄文第二版里是整整一章，第三版压成一段，只留结论与一处反证。

第三处是两处补写。第三章开头补出洪武元年那批作坊的来历，第四章开头补出太医院那一脉的来路。

第四处是新增的三段。第一章补出四次重启的证据来源：它们没有文献依据，只有岩芯与冰芯。第二章在会昌五年之后补一段，说本书所据的卷宗里也有几处同样的空白。第三章在宝船档案之后补一段，说那次消失比焚毁所能解释的更彻底。

这三段与两处补写不属于勘误，是新写的。写的时候有一种奇怪的感觉：每补一段，能引的材料就少一分。

一本书写了几十年，作者与它的关系是会变的。最初是他在写它，后来是他在读它，最后读它的时候，与读一个陌生人的书差不多。

\vspace{2mm}
\begin{flushright}
{\small А.\,В.\,巴兰尼科夫}\quad{\footnotesize\color{pencil} 二〇二四年一月，于莫斯科}
\end{flushright}

\tocandfigs

